\section{Finite benchmarks for everyday causal cognition}
\label{annex:p19-rat72-benchmarks}

\begin{annexitem}{annex:p19.rat72-power-benchmark}{thm:p19-rat72-power-benchmark}{Independent generators and the ceiling boundary}
\begin{proof}
Conditioning on \(C=0\) makes \(E=K\). Independence of \(K\) and
\(C\) therefore gives \(q_0=r\). Conditioning on \(C=1\) makes
\(E=K\lor T\). By mutual independence, the probability that both
generators fail, conditional on \(C=1\), is \((1-r)(1-s)\).
Consequently \(q_1=1-(1-r)(1-s)=r+(1-r)s\). Subtraction gives
\(\Delta P=(1-r)s\). Positive probabilities for both values of \(C\)
ensure the two conditional probabilities are defined.

When \(r<1\), division by the strictly positive \(1-r=1-q_0\)
uniquely gives \(s\). When \(r=1\), \(K=1\) almost surely and
\(E=1\) almost surely under either observed value of \(C\), for
any \(s\). Distinct choices \(s=0\) and \(s=1\), for example,
produce the same two conditional effect probabilities, proving the
stated nonidentification. With two independent generators of success
probability \(0.75\), the same complement calculation gives
\(1-0.25^2=0.9375\); no addition of probabilities of overlapping
success events is licensed.
\end{proof}
\end{annexitem}

\begin{annexitem}{annex:p19.rat72-diamond-benchmark}{thm:p19-rat72-diamond-benchmark}{The complete observational and intervention calculation}
\begin{proof}
Each supplied conditional factor is a normalized binary mass function.
Summing the product first over \(d\), then over \(b,c\), then over
\(a\), gives one, so the stipulated factorization defines a joint law.
Summing over \(b,d\) gives
\(P(A=a,C=c)=P(A=a)\ell_a(c)\). The assumed positive marginal for
\(C=c\) permits conditioning, giving \(P(A=1\mid C=c)=k_c\)
and \(P(A=0\mid C=c)=1-k_c\). Summing the conditional joint law
over \(a\) shows
\[
 P(B=1\mid C=c)=\sum_{a=0}^1P(A=a\mid C=c)P(B=1\mid A=a)
 =k_cu_1+(1-k_c)u_0=h_c.
\]
The factorization also gives
\(P(D=1\mid B=b,C=c)=t_{bc}\) for positive-probability conditioning
cells. Zero-mass cells contribute zero to the sum and can retain their
supplied kernels. Thus total probability yields
\[
 P(D=1\mid C=c)
 =\sum_{a=0}^1\sum_{b=0}^1
 P(A=a\mid C=c)P(B=b\mid A=a)t_{bc}
 =t_{1c}h_c+t_{0c}(1-h_c).
\]
This writes separate correctly indexed sums for \(c=0\) and \(c=1\).

Fixed replacement of \(C\)'s mechanism sets \(C=c\) and deletes
only its incoming arrow. Under the benchmark's unchanged independent
mechanisms, the remaining joint law is
\[
 P^{\operatorname{do}(C=c)}(a,b,d)=P(a)P(b\mid a)P(d\mid b,c).
\]
Its normalization follows by summing over \(d,b,a\). Its marginal
for \(A\) is therefore unchanged, and its marginal for \(B=1\) is
\(ru_1+(1-r)u_0=h\). Summing over \(a,b\) now gives
\[
 P(D=1\mid\operatorname{do}(C=c))
 =\sum_{a=0}^1\sum_{b=0}^1P(A=a)P(B=b\mid A=a)t_{bc}
 =t_{1c}h+t_{0c}(1-h).
\]
These calculations give A2--A3 and A6--A8 with binary values stated
explicitly; none requires conditioning on an event of positive mass
in the intervention model.

For the printed finite calibration, substitution gives
\(P(C=1)=0.475\), \(P(C=0)=0.525\),
\(k_1=0.45/0.475=18/19\) and
\(k_0=0.05/0.525=2/21\). Hence
\(h_1=0.1+0.75(18/19)=77/95\),
\(h_0=0.1+0.75(2/21)=6/35\), and \(h=0.475\).
The present-\(C\) predictive probabilities are
\(0.75+0.18(77/95)\) and \(0.75+0.18(0.475)=0.8355\).
The absent-\(C\) probabilities are
\(0.75(6/35)=9/70\) and \(0.75(0.475)=0.35625\).
Multiplication by forty and rounding to the nearest integer gives
the source's diagnostic counts \(38,4,20,20\) and predictive counts
\(36,5,33,14\). Figure~1's printed \(0.93\) is the value used here;
the two independent \(0.75\) generators would instead give the exact
\(0.9375\), without changing these rounded counts. This distinction
does not alter any recorded participant mean or SD.
\end{proof}
\end{annexitem}
